\chapter{Практика}

\section{Подготовка к построению модели}

Перед построением модели построим попарные графики зависимости признаков друг от друга с помощью библиотеки seaborn

\begin{lstlisting}
    sns.pairplot(
    df[pairplot_features + ["landmass"]],
    hue="landmass",
    diag_kind="hist"
)
\end{lstlisting}

\begin{figure}[H]
    \centering
    \includegraphics{10.png}
    \caption{Попарная зависимость признаков}
\end{figure}

По графикам можно сказать, что признаки не обладают хорошей "отделяемостью", что очень скажется на качестве модели.


\section{Построение модели}

Построим модель несколько раз, но с разным количеством соседей, чтобы определить оптимальное кол-во соседей. На каждой итерации будем
высчитывать accuracy модели и сопоставлять с k

\begin{lstlisting}
for k in k_values:
    model = Pipeline(steps=[
        ("preprocessor", preprocessor),
        ("classifier", KNeighborsClassifier(n_neighbors=k))
    ])
    
    model.fit(X_train, y_train)
    
    y_train_pred = model.predict(X_train)
    y_test_pred = model.predict(X_test)
    
    train_scores.append(accuracy_score(y_train, y_train_pred))
    test_scores.append(accuracy_score(y_test, y_test_pred))
\end{lstlisting}

Получившийся график:


\begin{figure}[H]
    \centering
    \includegraphics{20.png}
    \caption{Зависимость точности модели от k}
\end{figure}

Точность модели колеблится около 0.5, что очень плохо. По сути модель просто на удачу подбрасывает монетку.
Оптимальное значение accuracy достигается при k=2, т.к при k=2 точность максимально отклоняется от 0.5 (хоть и в отрицательную сторону). \\

Ошибка cross-validation:

\begin{figure}[H]
    \centering
    \includegraphics{30.png}
    \caption{Зависимость ошибки cross-validation от k}
\end{figure}

Теперь построим модель с k=2 и выведем report по модели:

\begin{figure}[H]
    \centering
    \includegraphics{40.png}
    \caption{Отчет по итоговой модели}
\end{figure}

Precision (точность) - оценка, показывающая, насколько можно доверять предсказаниям этого класса. Вычисляется как $\frac{TP}{TP+FP}$.\\

Recall (полнота), показывает, насколько хорошо модель находит этот класс. Вычисляется как $\frac{TP}{TP+FN}$.\\

F1 оценка - показывает общую адекватность модели по классу. Вычисляется как $2 * \frac{precision*recall}{precision+recall}$. \\

macro avg - среднее по классам.\\

weight avg - среднее с учётом количества объектов.












\endinput